%! Factoring Special Forms — Unit 1, Lesson 1.5 — Exit Ticket (Answer Key)
\documentclass[10pt]{article}
\usepackage{atda-article}
\usepackage{atda-key}   % pulls in -boxes; do NOT also load -boxes
\newcommand{\TallMath}[1]{$\displaystyle #1\rule[-1.4em]{0pt}{3.2em}$}
\setlength{\workrowsep}{12pt}

\begin{document}

\pageheader{Unit 1, Lesson 1.5}{Exit Ticket --- Answer Key}
% NAMESTRIP: no \namedateperiod here — mirrors the blank. Teacher-only prose goes
% in the lesson plan as \begin{teachernote}[Exit Ticket], never in a key.

\vspace{0.15in}

\begin{notesbox}{Special Forms --- On Your Own}
\small
Notes closed. Three items. Name the pattern, factor completely, and check by multiplying back.

\begin{enumerate}[leftmargin=1.6em, itemsep=8pt, topsep=6pt]

\item Factor: $x^2-14x+49$
\begin{work}
  2(x)(7) &= 14x \quad\checkmark \\
  x^2-14x+49 &= (x-7)^2
\end{work}

\item Factor: $x^3-64$
\begin{work}
  x^3-64 &= x^3-4^3 \\
         &= (x-4)\left(x^2+4x+16\right)
\end{work}

\item A student factors $2x^2-18$ as $(2x+6)(x-3)$. Is the answer correct? Is it complete?
Explain in one sentence, then give the complete factorization.
\begin{work}
  2x^2-18 &= 2\left(x^2-9\right) \\
          &= 2(x+3)(x-3)
\end{work}
\ansline{Correct but not complete --- $2x+6$ still shares a $2$; the GCF belonged out front first.}

\end{enumerate}
\end{notesbox}

\end{document}
